% APPENDIX B: A Complete Runnable Fourier Neural Operator
% k_max UNIFIED WITH PART I (Ch. 1.5, eq:ch05-modeset, and Ch. 1.6): k_max is the
% highest retained wavenumber and equals 11 here; the script's n_modes = 12 is a mode
% count -- twelve modes per axis, indices 0,...,11. Listings are byte-frozen and still
% carry the script's own variable names; only prose, the table and its caption changed.
% Typesetting pass: Table B.1 was set as a breakable xltabular in the text flow rather than a
% float, because as a [htbp] tabularx it owned a page to itself; its third column is
% raggedright, which is also what removes the underfull cells.
% Back matter, Appendix B: the book's worked example, whose reference implementation is
% embedded verbatim as Listings~\ref{lst:appb-data}--\ref{lst:appb-eval}, together with
% the caveats stated with the experiment.
% Code line numbers follow the reference file: the four listings carry its 166 lines,
% split BEFORE lines 60, 105 and 140, those three blank separators opening the listing
% that follows each split. The listings environment drops an empty last line instead of
% numbering it, so the previous arrangement (separators closing the preceding block) left
% the gutter numbers skipping 60, 105 and 140. No line is omitted, altered,
% reformatted, or reordered.

% The reference implementation is printed verbatim, which requires two departures
% from the shared pystyle, both local to this appendix and neither a change to the
% code: the escape character is disabled, because line 148 of the script contains a
% literal backslash, and the em dash on line 2 is given a literate replacement so
% that listings typesets the UTF-8 byte sequence instead of choking on it.
\lstdefinestyle{appbpystyle}{
    style=pystyle,
    escapechar={},
    literate={—}{{\textemdash}}1
}

\chapter{A Complete Runnable Fourier Neural Operator}\label{app:b}

Part I proves that a spectral network with fixed width and a fixed retained band
defines a map on functions; Part II builds such networks at two and three spatial
dimensions, with production layouts, schedulers, and export contracts. Neither
supplies an instance small enough to read at one sitting and cheap enough to run
without a graphics processing unit (GPU). This appendix is that instance: the
one-dimensional periodic heat-equation Fourier neural operator (FNO)
\cite{li2021fourier} that the book takes as its worked example, with its
reference implementation \texttt{FNO Worked Example (PyTorch).py} printed here in
full and unaltered. It makes three things inspectable at once: the analytic target
Part I's error figures are measured against, the configuration whose mode budget
\chapref{ch:part1-the-fourier-neural-operator-spectral-theor}
reasons about, and the pair of numbers drawn as markers in
Figure~\ref{fig:super_resolution}.

Everything below consolidates. The input family, the hyperparameters, the measured
errors, and the qualifications attached to them come from the reference script and from
the record of the run it documents, not from any new computation; where that record is
silent --- timings, hardware beyond
the statement that a central processing unit (CPU) suffices, repeated-seed
variability, any resolution other than the two reported --- the text says so rather
than filling the gap. Prerequisites are Fourier series on $[0,1)$, the spectral layer
\eqref{eq:ch23-layer}, the relative $L^2$ metric \eqref{eq:rel_l2_loss}, and the
PyTorch module interface. What follows is the problem and its exact solution, the
configuration, the script, the measurement, and the limits of the measurement.

\section{The Problem, Its Exact Solution, and the Input Measure}
\label{sec:appb-problem}

The model problem is the periodic heat equation of the worked example,
\begin{equation}
  \label{eq:appb-pde}
  \partial_t u = \alpha\,\partial_{xx} u, \qquad u(x,0) = u_0(x), \qquad
  x \in [0,1) \ \text{periodic}, \qquad \alpha = 0.01,\quad T = 1.0,
\end{equation}
whose constants the script states in its opening docstring. Since
$\partial_{xx} e^{2\pi i k x} = -(2\pi k)^2 e^{2\pi i k x}$, the coefficients of a
periodic solution obey $\partial_t \hat{u}(k,t) = -\alpha (2\pi k)^2 \hat{u}(k,t)$:
each mode decays alone, and the solution operator is a Fourier multiplier,
\begin{equation}
  \label{eq:appb-multiplier}
  u(x,T) = \mathcal{F}^{-1}\!\left[ \hat{u}_0(k)\, e^{-\alpha (2\pi k)^2 T} \right](x).
\end{equation}
That is why no numerical solver appears anywhere in the script: a forward
real-to-complex fast Fourier transform (FFT), one multiplication per mode, and an
inverse transform give the target field exactly, up to numerical precision, on any
grid resolving the input. The same factor appears in Figure~\ref{fig:spectral_decay},
and \chapref{ch:part1-the-fourier-neural-operator-spectral-theor} records that it
reaches $1.1\times10^{-11}$ by $k=8$ at these values of $\alpha$ and $T$: truncation is
negligible here, so the observed error is a network-capacity error. The operator to be
learned is $\mathcal{G}\colon u_0(\cdot) \mapsto u(\cdot,T)$, a linear map.

The input measure is fully specified:
\begin{equation}
  \label{eq:appb-input}
  u_0(x) = \sum_{k=1}^{8} c_k\, k^{-3/2} \cos(2\pi k x + \phi_k), \qquad
  c_k \sim \mathcal{N}(0,1), \quad \phi_k \sim U(0,2\pi),
\end{equation}
independent standard normal amplitudes, phases uniform on $[0,2\pi)$, a fixed
$k^{-3/2}$ envelope, wavenumbers $1$ through $8$. Three hundred functions are drawn
for training, fifty for validation, and fifty fresh draws of the same family for
testing --- new samples of \eqref{eq:appb-input}, not a new kind of input. Because the
band stops at eight, every input is resolved with room to spare on the sixty-four-point
training grid, whose real transform carries $F=\lfloor 64/2\rfloor+1=33$ coefficients.
That fact governs Section~\ref{sec:appb-scope}: it is why the fine-grid evaluation can
transfer at all, and why it cannot be read as recovering what the coarse grid lost.

\section{Configuration and Architecture}
\label{sec:appb-config}

Table~\ref{tab:appb-config} collects every setting stated for this experiment,
with the script line that fixes it; its first two rows are the constants of
\eqref{eq:appb-pde}. The
model is the lift--layer--project composition of \cite{kovachki2023neural}, one
Fourier layer between a pointwise lift and a pointwise projection,
\begin{equation}
  \label{eq:appb-layer}
  \text{out} = \mathcal{Q}\!\left[ \tanh\!\left( \mathcal{W}_1 v_1
  + \mathcal{F}_N^{-1}\!\left( R_\phi \hat{v}_1 \right) + b_1 \right) \right],
  \qquad v_1 = \mathcal{P}u, \quad \hat{v}_1 = \mathcal{F}_N v_1,
\end{equation}
which is \eqref{eq:ch23-layer} at depth $L=1$, one channel in and out, and width
$d_{\text{hidden}}=16$, except that the bias $b_1$ is written out, as the script has
it and \eqref{eq:ch23-layer} does not. Two conventions need care. The first is the
book's: $k_{\max}$ denotes the highest retained wavenumber, as in the mode set
\eqref{eq:ch05-modeset}, and never the number of retained coefficients. The script's
parameter is that number instead, so \texttt{n\_modes}=12 retains twelve modes per
axis, indices $0,\dots,11$, whence $k_{\max}=11$ and not $12$ --- the distinction is
flagged because the constant mode is counted --- and the script's
\texttt{min(n\_modes, F)}
keeps the slice inside the half-spectrum the grid supplies. The spectral branch is two
real tensors of shape $12\times16\times16$
rather than one complex array: the real and imaginary parts of a single band of channel
matrices. The two-dimensional layer of \eqref{eq:ch23-Rtensors} holds two \emph{complex}
tensors instead, because its retained set \eqref{eq:ch23-corners} has an upper and a
lower corner, and the tally \eqref{eq:ch23-paramcount} counts both blocks; a
one-dimensional half-spectrum supplies only the single band this script stores, so
\eqref{eq:ch23-paramcount} has no reading here.

The differences from Part II's reference implementation are conventional, not
substantive: the channel axis last rather than the channel-first layout
\eqref{eq:ch21-grid-layout}; \texttt{einsum} rather than $1\times1$ convolutions;
$\tanh$ where Listing~\ref{lst:ch23-fno2d} uses the Gaussian error linear unit (GELU);
one projection matrix rather than a two-layer network; no appended coordinate
channels; and the complex product expanded by hand into real and imaginary parts and
written into a zero-filled array of the full transform length, the allocation-as-padding
device the two-dimensional convolution of
\chapref{ch:part2-fourier-neural-operators-fno-from-scratch} also uses. Where
\cite{li2021fourier} imposes conjugate symmetry on $R_\phi$ because the kernel is
real, \texttt{irfft} here reconstructs from a half-spectrum, with no such constraint
in the parameters. The conclusion Table~\ref{tab:appb-config} carries is the one
\chapref{ch:part1-the-fourier-neural-operator-spectral-theor} argues for: the retained
band $0\le k\le k_{\max}=11$, twelve coefficients, covers every wavenumber
\eqref{eq:appb-input} contains, so the error of Section~\ref{sec:appb-result} is not
truncation.

{\small
\begin{xltabular}{\textwidth}{@{}l l >{\raggedright\arraybackslash}X@{}}
\caption{\textbf{The configuration of the worked example as stated, with the script line
that fixes each value.} Line numbers refer to \texttt{FNO Worked Example (PyTorch).py}.
The band entry clarifies a counting convention: the script
fixes \texttt{n\_modes}=12, a mode \emph{count}, which retains twelve modes per axis,
indices $0,\dots,11$; in the book's convention of \eqref{eq:ch05-modeset}, the highest
retained wavenumber is therefore $k_{\max}=11$, not $12$. Nothing in this table is
chosen here; every row is a quotation of a number that appears in the code or in the
record kept with it.}
\label{tab:appb-config}\\
\toprule
\textbf{Item} & \textbf{Value} & \textbf{Where it is fixed} \\
\midrule
\endfirsthead
\multicolumn{3}{r}{\textit{Table~\ref{tab:appb-config}, continued}}\\
\toprule
\textbf{Item} & \textbf{Value} & \textbf{Where it is fixed} \\
\midrule
\endhead
\midrule
\multicolumn{3}{r}{\textit{Table~\ref{tab:appb-config}, continued on the next page}}\\
\endfoot
\bottomrule
\endlastfoot
Diffusivity $\alpha$ & $0.01$ & line 22; docstring line 6 \\
Horizon $T$ & $1.0$ & line 22; docstring line 6 \\
Input band & wavenumbers $1,\dots,8$, envelope $k^{-3/2}$ & \texttt{n\_modes=8} default, lines 28--44 \\
Ensemble & $300$ train, $50$ validation, $50$ held-out & lines 110--111, 144 \\
Training grid $N$ & $64$, at $x_i=i/N$ & lines 109, 56 \\
Evaluation grid & $256$, same functions, no retraining & line 153 \\
Retained band $k_{\max}$ & $11$ & line 73: \texttt{n\_modes=12}, i.e.\ twelve modes per axis, indices $0,\dots,11$; used at line 91 \\
Hidden width $d_{\text{hidden}}$ & $16$ & line 73 \\
Channels $C_{\text{in}},C_{\text{out}}$ & $1$, $1$ & line 73 \\
Fourier layers $L$ & $1$ & class \texttt{TinyFNO}, lines 64--104 \\
Activation & $\tanh$ & line 103 \\
Precision & \texttt{float32} & line 23 \\
Seeds & \texttt{torch.manual\_seed(0)}, \texttt{np.random.seed(0)} & lines 20--21 \\
Optimizer & Adam, learning rate $10^{-3}$ & line 116 \\
Schedule & step decay, halve every $100$ epochs & line 117 \\
Epochs, batch size & $400$, $32$ & line 124 \\
Training loss & mean of squared per-sample relative $L^2$ & lines 131--133 \\
Reported metric & mean of per-sample relative $L^2$ & \texttt{rel\_l2}, lines 119--122 \\
\end{xltabular}
}

\section{The Script}
\label{sec:appb-code}

The implementation follows in four blocks --- data, model, training, evaluation ---
printed byte-for-byte from the reference file and numbered as the file numbers it, so
the same content sits on the same lines in both. All 166 lines are present, and the
gutter numbers run unbroken from 1 to 166: the blocks break \emph{before} lines 60,
105 and 140, so each of those three blank separator lines opens the listing that
follows the split rather than closing the one before it; nothing is reformatted,
renamed, reordered, or improved.
Running the file trains on a CPU and writes an archive \texttt{fno\_results.npz} that
the companion plotting script reads.

\subsection{Data generation and exact targets}

\begin{lstlisting}[style=appbpystyle, firstnumber=1, caption={\texttt{FNO Worked Example (PyTorch).py}, lines 1--59: module docstring, seeds, the input family \eqref{eq:appb-input}, the exact periodic heat solution \eqref{eq:appb-multiplier}, and the grid sampler. Printed verbatim.}, label={lst:appb-data}]
"""
fno_example.py — Fourier Neural Operator tutorial example (PyTorch, CPU).

Problem: 1-D periodic heat diffusion.
    u_t = alpha * u_xx,   u(x,0) = u_0(x),   x in [0,1) periodic
    alpha = 0.01, T = 1.0.

Analytic solution (spectral; exact, up to numerical precision, on grids that resolve these band-limited periodic inputs):
    u(x,T) = F^{-1}( F(u_0)_k * exp(-alpha (2 pi k)^2 T) )

Target operator:  G: u_0(.) -> u(.,T)    (function-to-function map)

We build a small FNO with one Fourier layer, train on random smooth
initial functions, then test on unseen initial functions AND on a finer
grid (zero-shot super-resolution: no retraining).
"""
import torch
import numpy as np

torch.manual_seed(0)
np.random.seed(0)
ALPHA, T_END = 0.01, 1.0
DTYPE = torch.float32

# ---------------------------------------------------------------
# 1) Data generation: random smooth periodic functions
# ---------------------------------------------------------------
def make_init_fns(n_funcs, n_modes=8):
    """Callables u0(x): random superpositions of cosines (smooth, periodic)."""
    funcs = []
    for _ in range(n_funcs):
        coeffs = np.random.randn(n_modes)
        phases = np.random.uniform(0, 2*np.pi, n_modes)
        amps   = 1.0 / np.arange(1, n_modes+1) ** 1.5
        funcs.append((coeffs, phases, amps))
    def make_u0(params):
        c, p, a = params
        def u0(x):
            s = np.zeros_like(x)
            for k in range(len(c)):
                s += c[k]*a[k]*np.cos(2*np.pi*(k+1)*x + p[k])
            return s
        return u0
    return [make_u0(p) for p in funcs]

def heat_solve_exact(u0_t, alpha=ALPHA, T=T_END):
    """Exact periodic heat solution. u0_t: torch (B,N,1)."""
    with torch.no_grad():
        N = u0_t.shape[1]
        c = torch.fft.rfft(u0_t, dim=1)                       # (B,F,1)
        k = torch.fft.rfftfreq(N, device=u0_t.device) * N
        decay = torch.exp(-alpha*(2*np.pi*k)**2 * T)[None, :, None]
        return torch.fft.irfft(c * decay, n=N, dim=1)          # (B,N,1)

def dataset(fns, N):
    x = np.linspace(0, 1, N, endpoint=False)
    U0 = np.stack([f(x) for f in fns])[:, :, None]
    U0 = torch.tensor(U0, dtype=DTYPE)
    return U0, heat_solve_exact(U0)
\end{lstlisting}

Two points in Listing~\ref{lst:appb-data} matter. \texttt{make\_init\_fns} returns
closures, not arrays, so one draw can be sampled on any grid --- the operation the
transfer claim depends on. \texttt{heat\_solve\_exact} is \eqref{eq:appb-multiplier}
transcribed: \texttt{rfft} along the point axis, \texttt{rfftfreq} scaled by $N$ to
integer wavenumbers, the decay broadcast over batch and channel, and \texttt{irfft}
back to physical space.

\subsection{The one-layer model}

\begin{lstlisting}[style=appbpystyle, firstnumber=60, caption={\texttt{FNO Worked Example (PyTorch).py}, lines 60--104: \texttt{TinyFNO}, the pointwise lift, the single spectral layer with the complex weight held as two real tensors, and the pointwise projection. Printed verbatim; the inline shape comments are the script's own; the block opens on the reference file's blank line 60.}, label={lst:appb-model}]

# ---------------------------------------------------------------
# 2) Tiny FNO in PyTorch (one Fourier layer)
# ---------------------------------------------------------------
class TinyFNO(torch.nn.Module):
    """
    v1  = P(u)                              lift (pointwise linear)
    v2  = tanh( W1 v1 + FFT^{-1}(R * FFT(v1)) + b1 )   Fourier layer
    out = Q(v2)                              project (pointwise linear)

    R: complex tensor (n_modes, d_hidden, d_hidden), applied per retained
    Fourier mode (channel mixing). Modes above n_modes are truncated.
    """
    def __init__(self, d_in=1, d_out=1, d_hidden=16, n_modes=12):
        super().__init__()
        self.P  = torch.nn.Parameter(torch.randn(d_in, d_hidden) / d_hidden**0.5)
        self.W1 = torch.nn.Parameter(torch.randn(d_hidden, d_hidden) / d_hidden**0.5)
        self.b1 = torch.nn.Parameter(torch.zeros(d_hidden))
        # complex weights: R[m, h, i]
        Rr = torch.randn(n_modes, d_hidden, d_hidden) / d_hidden**0.5
        Ri = torch.randn(n_modes, d_hidden, d_hidden) / d_hidden**0.5
        self.Rr = torch.nn.Parameter(Rr)
        self.Ri = torch.nn.Parameter(Ri)
        self.Q  = torch.nn.Parameter(torch.randn(d_hidden, d_out) / d_hidden**0.5)
        self.n_modes = n_modes

    def forward(self, u):
        B, N, _ = u.shape
        v1 = torch.einsum('bnm,mh->bnh', u, self.P)            # (B,N,H)
        v1_f = torch.fft.rfft(v1, dim=1)                        # (B,F,H) complex
        F = v1_f.shape[1]
        m = min(self.n_modes, F)
        vf = v1_f[:, :m]                                         # (B,m,H)
        vr, vi = vf.real, vf.imag
        # complex multiply as real arithmetic:
        # (vr + i vi) * (Rr + i Ri) = (vr Rr - vi Ri) + i (vr Ri + vi Rr)
        Rf_re = torch.einsum('bfh,fhi->bfi', vr, self.Rr[:m]) - torch.einsum('bfh,fhi->bfi', vi, self.Ri[:m])
        Rf_im = torch.einsum('bfh,fhi->bfi', vr, self.Ri[:m]) + torch.einsum('bfh,fhi->bfi', vi, self.Rr[:m])
        Rf = torch.zeros(B, F, v1.shape[2], dtype=torch.complex64,
                         device=u.device)
        Rf[:, :m] = Rf_re.to(torch.complex64) + 1j*Rf_im.to(torch.complex64)
        ker = torch.fft.irfft(Rf, n=N, dim=1)                    # (B,N,H)
        z = torch.einsum('bnh,hi->bni', v1, self.W1) + ker + self.b1
        a = torch.tanh(z)
        return torch.einsum('bnh,hk->bnk', a, self.Q)
\end{lstlisting}

Listing~\ref{lst:appb-model} is \eqref{eq:appb-layer} with nothing hidden. The lift
and the local path act pointwise through \texttt{einsum}; the band is sliced off the
half-spectrum, multiplied mode by mode, written into a zero-filled array of length
$F$, and inverted, so truncation is literally the vanishing of the coefficients above
index $11$, that is, above $k_{\max}=11$. Both grids satisfy $F>k_{\max}$, since $F$ is
$33$ at $N=64$ and $129$ at $N=256$, so the band is twelve modes wide at $N=64$ and at
$N=256$ alike --- the sense in which one parameter set denotes one spectral projection
on both.

\subsection{Training}

\begin{lstlisting}[style=appbpystyle, firstnumber=105, caption={\texttt{FNO Worked Example (PyTorch).py}, lines 105--139: the $N=64$ ensembles, the model instance with the hyperparameters of Table~\ref{tab:appb-config}, the reporting metric \texttt{rel\_l2}, and the training loop with its squared relative-$L^2$ objective \eqref{eq:appb-loss}. Printed verbatim; the block opens on the reference file's blank line 105.}, label={lst:appb-train}]

# ---------------------------------------------------------------
# 3) Train
# ---------------------------------------------------------------
N = 64
train_fns = make_init_fns(300)
val_fns   = make_init_fns(50)
U0_tr, UT_tr = dataset(train_fns, N)
U0_va, UT_va = dataset(val_fns, N)

model = TinyFNO(d_in=1, d_out=1, d_hidden=16, n_modes=12)
opt = torch.optim.Adam(model.parameters(), lr=1e-3)
sched = torch.optim.lr_scheduler.StepLR(opt, step_size=100, gamma=0.5)

def rel_l2(pred, true):
    num = (pred - true).norm(dim=(1,2))
    den = true.norm(dim=(1,2)).clamp_min(1e-9)
    return (num/den).mean().item()

EPOCHS, BATCH = 400, 32
for epoch in range(1, EPOCHS+1):
    perm = torch.randperm(len(U0_tr))
    for b in range(0, len(U0_tr), BATCH):
        sel = perm[b:b+BATCH]
        out = model(U0_tr[sel])
        # per-sample relative L2 squared (correct formula, no extra 1/N)
        num = (out - UT_tr[sel]).square().sum(dim=(1,2))
        den = UT_tr[sel].square().sum(dim=(1,2)).clamp_min(1e-12)
        loss = (num / den).mean()
        opt.zero_grad(); loss.backward(); opt.step()
    sched.step()
    if epoch % 50 == 0:
        with torch.no_grad():
            val_out = model(U0_va)
        print(f"epoch {epoch:3d}  val rel-L2 {rel_l2(val_out, UT_va):.4f}")
\end{lstlisting}

Nothing in the optimization changes because inputs and outputs are functions; only the
architecture of the forward map does. The objective at lines
131--133 of Listing~\ref{lst:appb-train} is the loss the example specifies,
\begin{equation}
  \label{eq:appb-loss}
  \mathcal{L} = \frac{1}{B}\sum_{i=1}^{B}
  \frac{\|\hat{u}^{(i)} - u^{(i)}\|_{L^2}^2}{\|u^{(i)}\|_{L^2}^2},
\end{equation}
the batch mean of the squared per-sample relative error, with no grid factor because
the weights of \eqref{eq:rel_l2_loss} cancel between numerator and denominator.
\eqref{eq:appb-loss} is what is minimized; the reporting function \texttt{rel\_l2}
uses the other convention, the root of the ratio
per sample and then the average;
\chapref{ch:part2-the-practitioner-s-paradigm-infinite-dimen} distinguishes the two
and notes that the percentages measured here are of the root-then-average form and are
in physical units, there being no normalizer in these four blocks.

\subsection{Held-out evaluation and zero-shot transfer}

\begin{lstlisting}[style=appbpystyle, firstnumber=140, caption={\texttt{FNO Worked Example (PyTorch).py}, lines 140--166: held-out evaluation at $N=64$, the zero-shot re-sampling of the same fifty functions at $N=256$ with unchanged weights, and the archive written for the plotting script. Printed verbatim; the block opens on the reference file's blank line 140; the two printed lines supply Table~\ref{tab:appb-measured}.}, label={lst:appb-eval}]

# ---------------------------------------------------------------
# 4) Evaluate on unseen initial functions (held-out)
# ---------------------------------------------------------------
test_fns = make_init_fns(50)
U0_te, UT_te = dataset(test_fns, N)
with torch.no_grad():
    pred = model(U0_te)
print(f"\n[64-pt grid]   held-out rel-L2 error: {rel_l2(pred, UT_te):.4f}")

# ---------------------------------------------------------------
# 5) Zero-shot super-resolution: same underlying functions, N=256
# ---------------------------------------------------------------
N2 = 256
U0_te2, UT_te2 = dataset(test_fns, N2)
with torch.no_grad():
    pred2 = model(U0_te2)
print(f"[256-pt grid]  held-out rel-L2 error (no retraining): {rel_l2(pred2, UT_te2):.4f}")

# save results for plotting
import os
np.savez(os.path.join(os.path.dirname(os.path.abspath(__file__)), 'fno_results.npz'),
         x=np.linspace(0,1,N,endpoint=False),
         u0=U0_te[0,:,0].numpy(), true=UT_te[0,:,0].numpy(), pred=pred[0,:,0].numpy(),
         x2=np.linspace(0,1,N2,endpoint=False),
         u0_2=U0_te2[0,:,0].numpy(), true2=UT_te2[0,:,0].numpy(), pred2=pred2[0,:,0].numpy())
print("saved fno_results.npz")
\end{lstlisting}

Everything in Listing~\ref{lst:appb-eval} after the training loop uses the weights as
they stand. The second evaluation differs from the first in one number, the count of
grid points the same functions are sampled at: the fifty test functions are drawn once
and resampled at both $64$ and $256$ points. That re-sampling is the whole of the
super-resolution experiment, and the archive written at the end holds its first sample
on each grid.

\section{The Measured Result}
\label{sec:appb-result}

Table~\ref{tab:appb-measured} reports what the run measured: observed numbers from the
account of one execution of the script above, not expectations derived from the
architecture, printed to four decimals by the two \texttt{print} calls of
Listing~\ref{lst:appb-eval}. They are transcribed, not reproduced --- the archive the
script writes is not among the artifacts held for this book, so no training was re-run
here. Validation error fell smoothly to roughly $1.6\%$ during
training; the held-out error was $1.41\%$ at $N=64$, and re-sampling those same
functions on a grid four times finer left the reported value unchanged at $1.41\%$. The
qualification that came with the number travels with it: both errors round to $0.0141$
at the printed precision, which is a statement about the report, not a claim that the two
underlying floating-point quantities were equal.

\begin{table}[htbp]
\centering
\caption{\textbf{Measured relative $L^2$ error on held-out inputs, as reported for
the run of the script in Section~\ref{sec:appb-code}.} The metric is
the root-then-average form of \eqref{eq:rel_l2_loss}, evaluated on the physical fields:
\texttt{rel\_l2} applies no pointwise normalizer, and the grid weights of
\eqref{eq:rel_l2_loss} cancel between numerator and denominator. The equality of the two
entries holds at the four decimals at which the script prints them. The record supplies
no timings, no spread across seeds, no per-epoch curve, and no error at any other
resolution.}
\label{tab:appb-measured}
\begin{tabular}{@{}llr@{}}
\toprule
\textbf{Evaluation} & \textbf{Retraining} & \textbf{Relative $L^2$ error} \\
\midrule
Held-out functions, $64$-point grid & none required (training grid) & $0.0141$ \ ($1.41\%$) \\
Same functions, $256$-point grid & none & $0.0141$ \ ($1.41\%$) \\
\bottomrule
\end{tabular}
\end{table}

What the table supports is transfer, not accuracy in the large: one parameter set,
fitted on three hundred functions sampled at sixty-four points, returned on fifty
held-out functions sampled at two hundred and fifty-six points the error it returned on
the same fifty at the training resolution --- no retraining, and no degradation at the
reported precision. The report adds that the pointwise error is small everywhere.
What the table does not support is comparison: $1.41\%$ is a modest error on an easy
target, and no competing network, classical solver, or third resolution appears alongside
the two reported values. The three-panel figure the plotting companion draws from the
archive is not among the book's supplied images, which is why the two values reach Part I
as the markers on the horizontal line
of Figure~\ref{fig:super_resolution}, whose caption repeats number and caveat alike.

\section{What the Experiment Establishes, and What It Does Not}
\label{sec:appb-scope}

Three qualifications come with this experiment, and they belong in the text, not a
footnote: an appendix is where a reader is likeliest to take the number at face value.

The first concerns the inputs. The functions of \eqref{eq:appb-input} carry eight
Fourier modes and the sixty-four-point grid resolves wavenumbers to $32$, so nothing
was lost between a function and its coarse samples: the $256$-point evaluation
resamples the \emph{same} already-resolved functions more finely. The experiment
demonstrates resolution transfer of the learned operator, not recovery of unresolved
high-frequency detail, and must not be cited as if it did --- the restriction the
caption of Figure~\ref{fig:super_resolution} places on the same data. What makes the
transfer legible is Proposition~\ref{prop:ch06-zeroshot} of
\chapref{ch:part1-discretization-invariance-and-continuum-li}, whose hypothesis
$N_{\text{train}} > 2k_{\max}$ holds here since $64 > 2\times 11 = 22$, and whose
band-limiting proviso holds too, the lift being pointwise linear and the family
\eqref{eq:appb-input} stopping at $k=8$, three wavenumbers below the cutoff
$k_{\max}=11$. That proposition speaks of one layer's spectral branch \emph{before} the
activation, which is exactly how far a single-layer model may go. The $\tanh$ at line 103
and the projection after it lie outside that identity.

The second concerns the equation. Heat diffusion is linear and
\eqref{eq:appb-multiplier} already solves it exactly and cheaply; applying the
multiplier to the samples is a spectral solve, which is how \texttt{heat\_solve\_exact}
makes the labels. This appendix draws the honest consequence --- this FNO is a teaching
example, not evidence that FNOs outperform classical solvers on this problem --- and
\chapref{ch:part1-the-fourier-neural-operator-spectral-theor} calls it a teaching
device for the same reason.

The third concerns reach. Both entries of Table~\ref{tab:appb-measured} are errors on
inputs drawn from \eqref{eq:appb-input}: smooth, band-limited, periodic, generated by
the training routine itself. The limitations recorded for this instance warn that
discretization invariance is not resolution-independent accuracy, that accuracy at an
unseen resolution must be checked empirically, and that an FNO trained on smooth initial
conditions may fail on discontinuous ones. The exercises proposed alongside it put those
questions to this script --- larger $\alpha$, \texttt{n\_modes} cut from $12$ to $4$, which is $k_{\max}$
cut from $11$ to $3$, evaluation at $32$ points, a second Fourier layer --- and none is
reported here, because none was measured. Surrogates earn their cost where a classical
solve is repeated many times
over a nonlinear, expensive, many-query problem:
\chapref{ch:part2-canonical-pde-benchmarks-and-reference-wor} treats this heat equation
as the preliminary to the Burgers, Darcy, and vorticity-form Navier--Stokes benchmarks
of Appendix~\ref{app:c}, and Appendix~\ref{app:d} collects the schedules and failure
modes a reader changing these settings will meet.

% ---------------------------------------------------------------------
\index{Fourier neural operator (FNO)!one-dimensional worked example}
\index{heat equation!runnable periodic example}
\index{band-limited input!eight-mode worked example}
\index{zero-shot super-resolution!measured at $N=64$ and $N=256$}
\index{relative $L^2$ error (metric)!root-then-average versus squared-ratio loss}
\index{discretization invariance!versus resolution-independent accuracy}
\index{spectral truncation!counting the constant mode in the band}
\index{learning rate scheduling!step decay in the worked example}
\index{hyperparameters!measured one-dimensional configuration}
\index{channel mixing!one-dimensional single-band instance}
